\textbf{Phần 1. Trắc nghiệm lựa chọn}
\setcounterpageref{ex}{0}

\begin{ex} %[5H5V29]
Một mảnh đất hình thang có đáy lớn $24\text{ m}$, đáy bé $16\text{ m}$. Nếu kéo dài đáy lớn thêm $5\text{ m}$ thì diện tích mảnh đất tăng thêm $35\text{ m}^2$. Diện tích mảnh đất ban đầu là
\choice
{$140\text{ m}^2$}
{\True $280\text{ m}^2$}
{$315\text{ m}^2$}
{$210\text{ m}^2$}
\loigiai{
Phần diện tích tăng thêm là một hình tam giác có đáy $5\text{ m}$ và cùng chiều cao với hình thang.\\
Chiều cao của mảnh đất là:
$35 \times 2 : 5 = 14\text{ (m)}$.\\
Diện tích mảnh đất ban đầu là:
$(24 + 16) \times 14 : 2 = 280\text{ (m}^2\text{)}$.\\
\textbf{Đáp án của câu hỏi là B}
}
\end{ex}

\begin{ex} %[5H5V29]
Một hình tròn có diện tích bằng $28{,}26\text{ cm}^2$. Chu vi của hình tròn đó là
\choice
{$9{,}42\text{ cm}$}
{\True $18{,}84\text{ cm}$}
{$28{,}26\text{ cm}$}
{$37{,}68\text{ cm}$}
\loigiai{
Tích hai bán kính của hình tròn là:
$r \times r = 28{,}26 : 3{,}14 = 9$.\\
Vì $3 \times 3 = 9$ nên bán kính hình tròn là $r = 3\text{ cm}$.\\
Chu vi của hình tròn là:
$2 \times 3{,}14 \times 3 = 18{,}84\text{ (cm)}$.\\
\textbf{Đáp án của câu hỏi là B}
}
\end{ex}

\begin{ex} %[5H5V29]
Cho tam giác $ABC$ có diện tích $120\text{ cm}^2$. Gọi $M, N$ lần lượt là trung điểm của hai cạnh $AB$ và $AC$. Diện tích của tứ giác $BMNC$ là
\choice
{$60\text{ cm}^2$}
{$80\text{ cm}^2$}
{\True $90\text{ cm}^2$}
{$100\text{ cm}^2$}
\loigiai{
Nối $B$ với $N$. Vì $N$ là trung điểm của $AC$ nên $S_{ABN} = \dfrac{1}{2} S_{ABC} = 60\text{ (cm}^2\text{)}$.\\
Vì $M$ là trung điểm của $AB$ nên $S_{AMN} = \dfrac{1}{2} S_{ABN} = 30\text{ (cm}^2\text{)}$.\\
Diện tích tứ giác $BMNC$ là:
$S_{BMNC} = S_{ABC} - S_{AMN} = 120 - 30 = 90\text{ (cm}^2\text{)}$.\\
\textbf{Đáp án của câu hỏi là C}
}
\end{ex}

\begin{ex} %[5H5V29]
Bánh xe của một xe tải có bán kính $0{,}4\text{ m}$. Khi bánh xe quay được đúng $500$ vòng thì xe tải đã đi được quãng đường dài bao nhiêu mét?
\choice
{$628\text{ m}$}
{\True $1\,256\text{ m}$}
{$2\,512\text{ m}$}
{$1\,000\text{ m}$}
\loigiai{
Chu vi của bánh xe là:
$2 \times 3{,}14 \times 0{,}4 = 2{,}512\text{ (m)}$.\\
Quãng đường xe tải đi được là:
$2{,}512 \times 500 = 1\,256\text{ (m)}$.\\
\textbf{Đáp án của câu hỏi là B}
}
\end{ex}

\begin{ex} %[5H5V29]
Một hình thang có đáy bé bằng $\dfrac{2}{3}$ đáy lớn, chiều cao bằng $12\text{ m}$ và diện tích bằng $300\text{ m}^2$. Độ dài đáy lớn của hình thang đó là
\choice
{$20\text{ m}$}
{$25\text{ m}$}
{\True $30\text{ m}$}
{$35\text{ m}$}
\loigiai{
Tổng độ dài hai đáy là:
$300 \times 2 : 12 = 50\text{ (m)}$.\\
Độ dài đáy lớn là:
$50 : (2 + 3) \times 3 = 30\text{ (m)}$.\\
\textbf{Đáp án của câu hỏi là C}
}
\end{ex}

\begin{ex} %[5H5V29]
Một cái sân hình chữ nhật dài $18\text{ m}$ và rộng $12\text{ m}$. Ở giữa sân người ta xây một bồn hoa hình tròn có bán kính $2\text{ m}$. Diện tích phần sân còn lại là
\choice
{$200{,}24\text{ m}^2$}
{\True $203{,}44\text{ m}^2$}
{$204{,}72\text{ m}^2$}
{$216\text{ m}^2$}
\loigiai{
Diện tích cái sân là:
$18 \times 12 = 216\text{ (m}^2\text{)}$.\\
Diện tích bồn hoa hình tròn là:
$3{,}14 \times 2 \times 2 = 12{,}56\text{ (m}^2\text{)}$.\\
Diện tích phần sân còn lại là:
$216 - 12{,}56 = 203{,}44\text{ (m}^2\text{)}$.\\
\textbf{Đáp án của câu hỏi là B}
}
\end{ex}

\begin{ex} %[5H5V29]
Một tam giác vuông có độ dài hai cạnh góc vuông là $15\text{ cm}$ và $20\text{ cm}$, cạnh huyền dài $25\text{ cm}$. Chiều cao tương ứng với cạnh huyền là
\choice
{$10\text{ cm}$}
{\True $12\text{ cm}$}
{$14\text{ cm}$}
{$16\text{ cm}$}
\loigiai{
Diện tích tam giác vuông là:
$15 \times 20 : 2 = 150\text{ (cm}^2\text{)}$.\\
Chiều cao tương ứng với cạnh huyền là:
$150 \times 2 : 25 = 12\text{ (cm)}$.\\
\textbf{Đáp án của câu hỏi là B}
}
\end{ex}

\begin{ex} %[5H5V29]
Hai hình tròn có tỉ số hai bán kính là $\dfrac{3}{4}$. Tỉ số diện tích của hình tròn bé và hình tròn lớn là
\choice
{$\dfrac{3}{4}$}
{$\dfrac{6}{8}$}
{\True $\dfrac{9}{16}$}
{$\dfrac{9}{12}$}
\loigiai{
Tỉ số diện tích hai hình tròn bằng bình phương tỉ số hai bán kính:
$$\left(\dfrac{3}{4}\right)^2 = \dfrac{9}{16}.$$
\textbf{Đáp án của câu hỏi là C}
}
\end{ex}

\begin{ex} %[5H5V29]
Một mảnh vườn hình thang vuông có đáy bé $10\text{ m}$, đáy lớn $15\text{ m}$ và chiều cao $8\text{ m}$. Người ta cắt bớt phần đất hình tam giác có đáy $5\text{ m}$ và chiều cao $8\text{ m}$ để làm lối đi. Phần đất còn lại là hình gì và có diện tích bằng bao nhiêu?
\choice
{Hình vuông, diện tích $64\text{ m}^2$}
{\True Hình chữ nhật, diện tích $80\text{ m}^2$}
{Hình chữ nhật, diện tích $100\text{ m}^2$}
{Hình thang, diện tích $80\text{ m}^2$}
\loigiai{
Khi cắt bỏ hình tam giác có đáy $15 - 10 = 5\text{ (m)}$ và chiều cao $8\text{ m}$, phần còn lại là hình chữ nhật có kích thước $10\text{ m} \times 8\text{ m}$.\\
Diện tích phần đất còn lại là:
$10 \times 8 = 80\text{ (m}^2\text{)}$.\\
\textbf{Đáp án của câu hỏi là B}
}
\end{ex}

\begin{ex} %[5H5V29]
Một hình tam giác có độ dài đáy tăng thêm $25\%$. Để diện tích tam giác không thay đổi thì chiều cao tương ứng phải
\choice
{Tăng thêm $25\%$}
{Giảm đi $25\%$}
{\True Giảm đi $20\%$}
{Giảm đi $15\%$}
\loigiai{
Khi đáy tăng thêm $25\%$ thì độ dài đáy mới bằng $125\% = 1{,}25$ đáy cũ.\\
Để diện tích không đổi thì chiều cao mới phải bằng:
$1 : 1{,}25 = 0{,}8 = 80\%$ chiều cao cũ.\\
Vậy chiều cao phải giảm đi:
$100\% - 80\% = 20\%$.\\
\textbf{Đáp án của câu hỏi là C}
}
\end{ex}

\begin{ex} %[5H5V29]
Một mặt bàn có hình dạng gồm một hình chữ nhật kích thước $1\text{ m} \times 2\text{ m}$ và ở hai đầu chiều rộng ghép thêm hai nửa hình tròn đường kính $1\text{ m}$. Diện tích của toàn bộ mặt bàn đó là
\choice
{$2{,}3925\text{ m}^2$}
{\True $2{,}785\text{ m}^2$}
{$3{,}14\text{ m}^2$}
{$2{,}57\text{ m}^2$}
\loigiai{
Diện tích phần hình chữ nhật là:
$1 \times 2 = 2\text{ (m}^2\text{)}$.\\
Hai nửa hình tròn ghép lại thành một hình tròn có đường kính $1\text{ m}$ (bán kính $r = 0{,}5\text{ m}$).\\
Diện tích hình tròn là:
$3{,}14 \times 0{,}5 \times 0{,}5 = 0{,}785\text{ (m}^2\text{)}$.\\
Diện tích toàn bộ mặt bàn là:
$2 + 0{,}785 = 2{,}785\text{ (m}^2\text{)}$.\\
\textbf{Đáp án của câu hỏi là B}
}
\end{ex}

\begin{ex} %[5H5V29]
Cho hình thang $ABCD$ có hai đáy $AB = 6\text{ cm}$ và $CD = 10\text{ cm}$. Biết diện tích tam giác $ABC$ bằng $24\text{ cm}^2$. Diện tích tam giác $ACD$ là
\choice
{$30\text{ cm}^2$}
{$36\text{ cm}^2$}
{\True $40\text{ cm}^2$}
{$48\text{ cm}^2$}
\loigiai{
Hai tam giác $ABC$ và $ACD$ có cùng chiều cao bằng chiều cao hình thang $ABCD$.\\
Do đó tỉ số diện tích bằng tỉ số độ dài hai đáy tương ứng:
$$\dfrac{S_{ACD}}{S_{ABC}} = \dfrac{CD}{AB} = \dfrac{10}{6} = \dfrac{5}{3}.$$
Diện tích tam giác $ACD$ là:
$24 \times \dfrac{5}{3} = 40\text{ (cm}^2\text{)}$.\\
\textbf{Đáp án của câu hỏi là C}
}
\end{ex}

\textbf{Phần 2. Trắc nghiệm đúng sai}
\setcounterpageref{ex}{0}

\begin{ex} %[5H5V29]
Một hình vuông có cạnh $10\text{ cm}$ vẽ tiếp xúc ngoài với một hình tròn bên trong (hình tròn có đường kính bằng cạnh hình vuông).
\choiceTF
{\True Bán kính của hình tròn là $5\text{ cm}$}
{\True Chu vi của hình tròn là $31{,}4\text{ cm}$}
{\True Diện tích của hình tròn là $78{,}5\text{ cm}^2$}
{Diện tích phần hình vuông nằm ngoài hình tròn là $25\text{ cm}^2$}
\loigiai{
\begin{itemchoice}
\item Đúng. Bán kính $r = 10 : 2 = 5\text{ (cm)}$.
\item Đúng. Chu vi $C = 2 \times 3{,}14 \times 5 = 31{,}4\text{ (cm)}$.
\item Đúng. Diện tích $S = 3{,}14 \times 5 \times 5 = 78{,}5\text{ (cm}^2\text{)}$.
\item Sai. Diện tích hình vuông là $10 \times 10 = 100\text{ (cm}^2\text{)}$. Diện tích phần nằm ngoài là: $100 - 78{,}5 = 21{,}5\text{ (cm}^2\text{)} 
eq 25\text{ cm}^2$.
\end{itemchoice}
}
\end{ex}

\begin{ex} %[5H5V29]
Một thửa ruộng hình thang có độ dài đáy lớn $36\text{ m}$, đáy bé bằng $\dfrac{2}{3}$ đáy lớn và chiều cao bằng trung bình cộng độ dài hai đáy.
\choiceTF
{\True Đáy bé của thửa ruộng là $24\text{ m}$}
{\True Chiều cao của thửa ruộng là $30\text{ m}$}
{\True Diện tích của thửa ruộng là $900\text{ m}^2$}
{Người ta dành $40\%$ diện tích thửa ruộng để trồng hoa, diện tích trồng hoa là $320\text{ m}^2$}
\loigiai{
\begin{itemchoice}
\item Đúng. Đáy bé: $36 \times \dfrac{2}{3} = 24\text{ (m)}$.
\item Đúng. Chiều cao: $(36 + 24) : 2 = 30\text{ (m)}$.
\item Đúng. Diện tích: $(36 + 24) \times 30 : 2 = 900\text{ (m}^2\text{)}$.
\item Sai. Diện tích trồng hoa là: $900 \times 40\% = 360\text{ (m}^2\text{)} 
eq 320\text{ m}^2$.
\end{itemchoice}
}
\end{ex}

\textbf{Phần 3. Trắc nghiệm trả lời ngắn}
\setcounterpageref{ex}{0}

\begin{ex} %[5H5V29]
Một hình tam giác có diện tích bằng $108\text{ cm}^2$ và cạnh đáy dài $18\text{ cm}$. Chiều cao tương ứng của hình tam giác đó là bao nhiêu xăng-ti-mét?
\shortans{12}
\loigiai{
Chiều cao của hình tam giác là:
$108 \times 2 : 18 = 12\text{ (cm)}$.
}
\end{ex}

\begin{ex} %[5H5V29]
Một mảnh đất hình thang có diện tích $210\text{ m}^2$, chiều cao $14\text{ m}$. Biết đáy lớn hơn đáy bé $6\text{ m}$. Tính độ dài đáy bé của mảnh đất theo đơn vị mét.
\shortans{12}
\loigiai{
Tổng độ dài hai đáy là:
$210 \times 2 : 14 = 30\text{ (m)}$.\\
Độ dài đáy bé là:
$(30 - 6) : 2 = 12\text{ (m)}$.
}
\end{ex}

\begin{ex} %[5H5V29]
Bán kính hình tròn thứ nhất là $6\text{ cm}$, bán kính hình tròn thứ hai là $2\text{ cm}$. Diện tích hình tròn thứ nhất gấp bao nhiêu lần diện tích hình tròn thứ hai?
\shortans{9}
\loigiai{
Tỉ số bán kính hai hình tròn là: $6 : 2 = 3$ (lần).\\
Tỉ số diện tích hai hình tròn là:
$3 \times 3 = 9\text{ (lần)}$.
}
\end{ex}

\begin{ex} %[5H5V29]
Một bánh xe đạp có bán kính $0{,}25\text{ m}$ lăn trên đường được quãng đường dài $314\text{ m}$. Hỏi bánh xe đạp đó đã quay được bao nhiêu vòng?
\shortans{200}
\loigiai{
Chu vi của bánh xe đạp là:
$2 \times 3{,}14 \times 0{,}25 = 1{,}57\text{ (m)}$.\\
Số vòng quay của bánh xe là:
$314 : 1{,}57 = 200\text{ (vòng)}$.
}
\end{ex}

\begin{ex} %[5H5V29]
Cho hình thang $ABCD$ có hai đáy $AB = 8\text{ cm}$ và $CD = 12\text{ cm}$, chiều cao hình thang là $9\text{ cm}$. Tính diện tích hình tam giác $BCD$ theo đơn vị xăng-ti-mét vuông.
\shortans{54}
\loigiai{
Hình tam giác $BCD$ có độ dài đáy $CD = 12\text{ cm}$ và chiều cao hạ từ đỉnh $B$ xuống đáy $CD$ chính bằng chiều cao hình thang là $9\text{ cm}$.\\
Diện tích tam giác $BCD$ là:
$12 \times 9 : 2 = 54\text{ (cm}^2\text{)}$.
}
\end{ex}

\begin{ex} %[5H5V29]
Một tấm thảm trải sàn hình chữ nhật có kích thước $4\text{ m} \times 6\text{ m}$. Ở chính giữa tấm thảm người ta thêu một hoa văn hình thoi có các đỉnh là trung điểm bốn cạnh của hình chữ nhật. Tính diện tích của hình thoi hoa văn đó theo đơn vị mét vuông.
\shortans{12}
\loigiai{
Hai đường chéo của hình thoi có độ dài bằng chính chiều dài và chiều rộng của tấm thảm hình chữ nhật ($6\text{ m}$ và $4\text{ m}$).\\
Diện tích của hình thoi là:
$6 \times 4 : 2 = 12\text{ (m}^2\text{)}$.
}
\end{ex}
